\section{从零开始的知识正文}
\label{sec:foundations}

\subsection{紧致性与连续性}
\label{sec:compactness}

\begin{definition}
设 \(X\) 为拓扑空间。若一族开集 \(\{U_i\}_{i\in I}\) 满足
\(X=\bigcup_{i\in I}U_i\)，就称它为 \(X\) 的一个\emph{开覆盖}。若每个开覆盖都能从中
选出有限多个仍覆盖 \(X\)，则称 \(X\) 为\emph{紧致空间}。
\end{definition}

紧致性的直觉是“无限的信息能由有限部分控制”。定义中的覆盖可以有任意多个成员，所以
证明紧致时不能只检查某一个熟悉覆盖；使用定理时也要确认空间、映射和拓扑确实满足假设。

\begin{theorem}[连续像保持紧致]
\label{thm:compact-image}
若 \(X\) 紧致，\(f:X\to Y\) 连续，则 \(f(X)\) 在子空间拓扑下紧致。
\end{theorem}

\begin{proof}
任取 \(f(X)\) 的开覆盖 \(\{V_i\}_{i\in I}\)。对每个 \(i\)，由子空间拓扑可取
\(Y\) 中开集 \(W_i\)，使 \(V_i=f(X)\cap W_i\)。连续性说明
\(f^{-1}(W_i)\) 在 \(X\) 中开。又因 \(\{V_i\}\) 覆盖 \(f(X)\)，任意 \(x\in X\)
的像 \(f(x)\) 属于某个 \(V_i\)，从而 \(x\in f^{-1}(W_i)\)。故
\(\{f^{-1}(W_i)\}_{i\in I}\) 是 \(X\) 的开覆盖。

由 \(X\) 紧致，存在 \(i_1,\dots,i_m\in I\)，使
\[
  X=\bigcup_{k=1}^{m}f^{-1}(W_{i_k}).
\]
任取 \(y\in f(X)\)，写成 \(y=f(x)\)。上式给出某个 \(k\) 使
\(x\in f^{-1}(W_{i_k})\)，故 \(y\in f(X)\cap W_{i_k}=V_{i_k}\)。因此
\(V_{i_1},\dots,V_{i_m}\) 覆盖 \(f(X)\)，这正是 \(f(X)\) 紧致的定义。
\end{proof}

证明的关键不是背诵符号，而是记住“把目标空间的覆盖沿 \(f\) 拉回定义域”。连续性保证
拉回后仍是开集，紧致性负责把无限覆盖压成有限覆盖，最后再把有限性推回像空间。

